\documentclass{article}

\PassOptionsToPackage{numbers}{natbib}
\usepackage[final]{neurips_2019}

\usepackage[utf8]{inputenc}
\usepackage[T1]{fontenc}
\usepackage{hyperref}
\usepackage{url}
\usepackage{booktabs}
\usepackage{amsfonts}
\usepackage{amsmath}
\usepackage{microtype}
\usepackage{graphicx}
\usepackage{xcolor}

\title{What Does a Wish Leave Out?\\[0.4em]
{\large A Measured Study of Requirement Recovery in Wish-Driven Code Agents}}

\author{
  Yanhao Bai \quad Kaitlyn Bassford \quad Wen Tao Koh \\
  Department of Computer Science, Columbia University \\
  \texttt{\{yb2630, kmb2297, wk2439\}@columbia.edu}
}

\raggedbottom

\begin{document}

\maketitle

\begin{abstract}
A wish-driven coding agent receives an intent and builds an artifact. When the wish leaves a
requirement implicit, the program can run, pass the tests the user remembered to mention, and
still violate what the user meant. We mechanically degrade benchmark specifications into
wishes by deleting four of the five requirement fields, train a small \emph{silence detector} to predict
which categories a wish leaves implicit, and use those predictions to target clarification
questions. This report measures the design on a 100-task development pilot under the
corrected protocol: the whole docstring prose before the examples is the intent,
HumanEval+/MBPP+ visible tests show inputs and observed output without expected values, and
BigCodeBench visible tests remain assertion-based. The run is nine conditions, three samples
per task, 2{,}700 units, one A100-SXM4-40GB session, 2.902 GPU-hours and \$0 paid API. The
specification gap is real: the full-specification ceiling B1 passes 62.3\% against 52.0\% for
the plain wish loop H0 (+10.3 percentage points, 95\% paired task-level bootstrap CI
[+3.0, +17.7]). Targeted clarification H3 is the only condition that recovers any of the gap
at the point estimate (53.7\%, recovery 0.16) and it beats generic questioning H2 (44.3\%)
by +9.3 pp [+1.7, +17.0]; but it does not beat the ceiling or sibling-function context H4
(51.7\%), and its difference from H0 (+1.7 pp [$-5.7$, $+9.0$]) crosses zero. Generic
questions hurt ($-7.7$ pp), withholding H3's answers hurts ($-11.0$ pp), acceptance tests
first is catastrophic ($-23.3$ pp), and the one-call baseline is weak (31.0\%). The detector
is a usable ranker but not a solved component (macro-F1 0.792; AUROC 0.952 constraints,
0.744 error behaviour, 0.632 examples). We state the corrected protocol, the per-category and
failure evidence, and the limits of a 100-task, single-session pilot.
\end{abstract}

\section{Introduction}

Wish-driven programming is attractive because the user states an outcome and the agent fills in
the rest. The same property makes failure quiet: an implementation can satisfy every stated
requirement and still miss a constraint, an error behaviour or a typing detail that the user
had in mind. Requirement elicitation is therefore part of the agent's job, not an optional
politeness. The design asks whether a small detector over the specification fields can make
that elicitation targeted: if the detector knows which category a wish omits, the agent can
ask about that category instead of spending a fixed question budget on generic prompts.

This report is the corrected measured version of that design. An earlier run produced a
null for targeted clarification, but its parser took only the first docstring sentence as the
intent, so a requirement stated later was filed as a constraint and deleted at L3, and its
HumanEval+/MBPP+ visible tests carried expected values, so the examples field was silently
restored. The committed levels use the corrected protocol: the intent is the whole prose block
before the examples, later unlabelled prose is part of the request, and HumanEval+/MBPP+
visible tests print observed output instead of expected values. BigCodeBench visible tests
stay assertion-based and are a stated residual. The same nine conditions were re-run on the
corrected levels; this report integrates that run and recomputes every headline number from
the committed unit traces.

The corrected answer is more useful than the earlier null. The specification gap is real and
about the same size: deleting the specification costs 10.3 points. Targeted clarification
now \emph{beats} generic clarification by 9.3 points, is the only condition that recovers any
of the gap at the point estimate, and does so while asking fewer questions and taking fewer
turns. It still does not reach the ceiling, and its gap over H0 and H4 crosses zero, so the
honest claim is relative: targeting is better than generic asking, not sufficient to close
the specification gap. Generic questions actively hurt, withholding the targeted answers
hurts more, and writing acceptance tests first from inputs-only visible tests is disastrous.
The detector ranks constraints well but its error-behaviour label remains weak.

The rest of the report gives the design and corrected protocol (Section~\ref{sec:approach}),
the measured run and its results (Section~\ref{sec:experiments}), the limitations that bound
the conclusion (Section~\ref{sec:limitations}), and the artifact map that makes each number
traceable (Section~\ref{sec:reproduce}).

\section{Related work}
\label{sec:related}

ClarifyGPT shows that intention clarification improves LLM code generation
\citep{mu2023clarifygpt}; our H2/H3 conditions are a controlled version of that idea, with the
clarification target chosen by a detector rather than by the model's own uncertainty. EvalPlus
supplies the rigorous hidden-test protocol we use for HumanEval+ and MBPP+
\citep{liu2023evalplus}, and BigCodeBench supplies the library-call stress case
\citep{zhuo2025bigcodebench}. Contemporary harnesses revise their own state online
\citep{xiao2026reharness}; we hold harness state fixed so that only the conditioning varies.
The planned H5 is offline harness evolution: RRSI regularizes recursive self-improvement
\citep{xia2026rrsi}, Harness$^2$ evolves harnesses without ground-truth verifiers
\citep{xu2026harness2}, and HarnessPro uses LLM-generated executable tests
\citep{pan2026harnesspro}. Those works bound what a future evolved harness could add; this
report is about a prompt-conditioned detector, not an evolved harness.

\section{Approach}
\label{sec:approach}

\subsection{Task, data and levels}

The agent receives a one-sentence wish and a Python file containing a function stub, and must
produce an implementation that passes a hidden unit-test suite. A unit is one function with a
shared stub, three visible tests and one hidden suite. We use HumanEval+ and MBPP+
\citep{liu2023evalplus}, which pair a docstring with one function, and BigCodeBench
\citep{zhuo2025bigcodebench}, which composes several library calls and is the stress case.
The source pool is 1{,}682 records; 133 were excluded (31 parse, 86 visible-test, 16
hidden-test failures), leaving 1{,}549 valid tasks, and the split keeps 1{,}100 (132
HumanEval+, 377 MBPP+, 591 BigCodeBench) as 500 detector-train / 200 evolution / 300
held-out plus a disjoint 100-task development pilot. The pilot's source mix is 12 HumanEval+, 34 MBPP+ and 54
BigCodeBench tasks. The 300 held-out tasks were never used for training, calibration,
threshold selection or the pilot.

Every specification is parsed mechanically into five fields: intent, signature and types,
examples, constraints and edge cases, and error behaviour. Levels are built by deleting
fields. L0 is the full specification; L1 drops examples and type detail; L2 additionally
drops constraints and error behaviour; L3 keeps only the intent prose. A field is never
invented, and an empty field stays empty. The hidden suite is byte-identical across levels, so
an L0--L3 difference is information loss rather than test variation. The labels for the
detector are field presence in L0: a positive means the category has separately rendered
content that the level construction deletes, though the whole-prose intent can still mention
some of it. Single-field drops L2$_c$ would attribute recovery to a category, but the
nine-condition pilot does not run them, so the per-category analysis below groups tasks by the
L0 field whose separately rendered content the wish omits rather than by an L2$_c$ contrast.

The corrected parser keeps the whole prose block before the examples as the intent. Later
unlabelled sentences are part of that request, not hidden constraints; explicit
\emph{Constraints}/\emph{Raises} labels and sentences carrying an error cue are additionally
classified into the constraint and error-behaviour fields, while the prose itself remains the
intent. This matters at L3: the corrected wish is a natural request rather than the first
sentence with later requirements deleted. The earlier levels used only the first sentence and
therefore attenuated both H0 and any asking effect; the corrected protocol fixes that.

L3 is smoothed by the generator before the pilot: 500 train and 100 pilot wishes were
rewritten into natural one-paragraph wishes with the same endpoint, and 12 of 600 fell back to
the raw intent. The detector is refit on the 500 smoothed train wishes. This keeps the wish
natural without changing the mechanical deletion protocol.

\subsection{Corrected visible tests}

Under the corrected protocol, HumanEval+ and MBPP+ visible tests exercise the same first three
inputs but print the observed output; expected values are removed from the source. The agent
therefore cannot recover the examples field from the visible suite, which is what the levels
protocol assumes. BigCodeBench's visible suite is a selection of \texttt{unittest} methods
whose expectations are assertions rather than literal \texttt{(inputs, results)} pairs, so it
remains assertion-based and still exposes worked examples in both protocols. This is a
residual: the examples field is not fully deleted for the BigCodeBench subset, and the
per-category examples row below is confounded by it. The same correction makes H1 weak by
construction: writing acceptance tests first is much less informative when the visible tests
contain no expected outputs, so H1's result is reported as measured with that caveat rather
than as evidence about acceptance-test practice in general.

\subsection{Harness and conditions}

The agent is a ReAct loop with exactly four tools: write the file, run the visible tests, read
the traceback of the last failed visible run, and ask the stakeholder one question. The loop
allows at most twelve assistant turns, has no repository browsing and no shell, and starts
from a fresh workspace per task. The stakeholder is a temperature-0 model that sees only the
L0 specification and answers ``That is not specified.'' when the answer is absent. This
follows the ClarifyGPT protocol of answering only from the specification
\citep{mu2023clarifygpt}.

All conditions share the tasks, tools, visible tests and turn cap. B1 gives the full L0
specification with no dialogue and serves as the ceiling. H0 is the plain wish loop. H1 may
write acceptance tests first, but only from the visible tests; under inputs-only tests it
receives no new information and very little expected behaviour. H2 asks generic stakeholder
questions through the tool within a five-question budget. H3 asks the detector-selected 2--3
questions up front, compiles the answers into a \emph{Requirements} block that conditions the
code prompt, and then runs one self-check-and-repair call. H4 adds five sibling functions with
doctests to the workspace. B2 draws five direct samples instead of asking. B3 calls the
endpoint once with only the wish and stub; with no separate frontier endpoint configured, it
uses the same Qwen3-8B at temperature 0.7. H3$^-$ asks H3's questions but withholds the
answers, separating asking from knowing. H5, offline harness evolution, remains a plan item
and is not in this pilot.

\subsection{Detector and question policy}

The detector is the frozen \texttt{Qwen/Qwen3-Embedding-0.6B} encoder (1024-dimensional, CPU)
plus a five-label linear head trained with binary cross-entropy and an L2 penalty. The
committed implementation uses 600 full-batch AdamW steps at learning rate 0.01 with L2 0.3;
inner folds select the penalty from $\{0.03, 0.1, 0.3\}$ and fit Platt calibration. The
deployed threshold is 0.5 for the reported detection metrics, while the H3 policy uses a
separately calibrated threshold of 0.1974 over calibrated presence times utility. Training
uses only the 500 train tasks; the 200 evolution, 300 held-out and 100 pilot tasks are
excluded from fitting, calibration, threshold selection and model selection. Intent and
signature/types are positive on all 500 train tasks, so their AUROC is undefined and the
detector is effectively a three-label head with 293 constraints, 491 examples and 85
error-behaviour positives.

The H3 policy assigns utilities (intent 0, examples 0.2, signature and types 0.5, constraints
1.0, error behaviour 1.0), ranks categories by calibrated presence times utility, and clamps
each task to two or three questions. Intent is never asked because it is already the wish. On
the deployed-style out-of-fold scores the policy asks 2.48 questions per task on average (260
tasks get two, 240 get three); informative-pair precision is 0.835, recall 0.554 and utility
recall 0.856. Selection rates are 0.70 for constraints, 0.37 for error behaviour, 0.42 for
examples, 1.00 for signature/types and 0 for intent. The pilot's H3 units average 2.48
questions, matching the policy.

\subsection{Metrics and paired bootstrap}

Pass@1 is the first sample of a unit; for B2 it is the first of the five internal direct
samples, and B2's pass rate is the fraction of units where any of the five passes. Recovery is
the fraction of the L0--L3 gap closed:
\[
  \mathrm{recovery}(c) = \frac{\mathrm{pass@1}(c) - \mathrm{pass@1}(H0)}
                                {\mathrm{pass@1}(B1) - \mathrm{pass@1}(H0)}.
\]
Per-category recovery groups tasks by the L0 category that the wish omits.

Contrasts are paired at the task level. Each task contributes one mean over its three samples
per condition; we resample the 100 tasks with replacement 2{,}000 times (seed
\texttt{20261121}), take the mean of the task-level differences in each resample, and report
the 2.5th and 97.5th percentiles. The same resample matrix is reused for every contrast. This
interval describes uncertainty about the average difference over tasks; it does not capture
variation across model seeds or sessions. The bootstrap uses \texttt{pass\_at\_1}, so B2's
contrast is defined on its first direct sample; B2's five-sample pass rate is reported
separately and its difference from H0 under that definition is $+2.7$ pp, which is not the
contrast in Table~\ref{tab:bootstrap}.

\section{Experiments}
\label{sec:experiments}

\subsection{Measured run}

The operator ran one A100-SXM4-40GB Colab session serving \texttt{Qwen/Qwen3-8B-AWQ} on vLLM
0.34.0 at temperature 0.7 and top-$p$ 0.95. The vLLM build did not expose native tool calls
(the endpoint rejected the native probe), so every unit used the JSON text protocol. All
2{,}700 units completed on the 100 development tasks with no held-out task. The run used
19{,}356{,}531 tokens (18{,}075{,}167 input, 1{,}281{,}364 output), 2.902 GPU-hours on two
config-stable shards (1.385 h and 1.517 h), and \$0 paid API. Failures are bucketed by the
exception the runner records, which says why a unit failed rather than only that it did.

\subsection{Main results}

Table~\ref{tab:conditions} gives all nine conditions. The specification gap is the largest
effect: B1 passes 62.3\% and H0 52.0\%, a 10.3-point difference whose interval excludes zero.
Targeted clarification is the only condition that recovers any of the gap at the point
estimate: H3 passes 53.7\%, recovery 0.162. It is also the only clarification condition that
beats another one: H3 $-$ H2 is $+9.3$ pp [$+1.7$, $+17.0$], while H2 is 7.7 points
\emph{below} H0 (44.3\%, [$-14.3$, $-1.0$]). The targeted condition does not reach the
ceiling and its difference from H0 is $+1.7$ pp [$-5.7$, $+9.0$], which crosses zero; H4 is
statistically indistinguishable from H3 (H3 $-$ H4 $+2.0$ pp [$-6.3$, $+9.7$]) and from H0
($-0.3$ pp [$-6.0$, $+5.3$]). H3$^-$ is 11.0 points below H3 (42.7\% against 53.7\%; difference $-11.0$ pp [$-18.7$, $-2.7$]),
so knowing the answers matters more than the act of asking them, but asking generically
appears to cost more than it gains. H1 is 23.3 points below H0 (28.7\%), a result that is
partly a protocol effect under inputs-only visible tests. B2's five samples raise its pass
rate to 54.7\% but its first-sample pass@1 is 51.3\% ($-0.7$ pp), so extra samples are not a
substitute for information. B3 is a weak one-call baseline at 31.0\%.

The corrected run thus supports a relative claim about targeting, not a claim that targeted
clarification closes the specification gap. H3's advantage over H2 is the first positive
clarification result in this study; it comes while H3 asks fewer questions than H2 (2.48 vs
2.94) and takes fewer turns (4.43 vs 6.36). The generic condition's extra interaction does not
buy accuracy, and its extra turns may actively cost. This is consistent with the targeted
condition converting each answer into a Requirements block and a single repair call, while H2
spends the budget on generic prompts that may not address the missing field.

\begin{table}[!ht]
  \centering
  \small
  \setlength{\tabcolsep}{4pt}
  \begin{tabular}{@{}llrrrrr@{}}
    \toprule
    Condition & Description & pass@1 & pass rate & recovery & questions & turns \\
    \midrule
    B1 & full L0 specification, no dialogue & 0.623 & 0.623 & 1.000 & 0.00 & 4.34 \\
    H0 & plain wish loop & 0.520 & 0.520 & 0.000 & 0.00 & 4.40 \\
    H1 & acceptance tests first & 0.287 & 0.287 & $-2.259$ & 0.00 & 4.51 \\
    H2 & generic clarification & 0.443 & 0.443 & $-0.743$ & 2.94 & 6.36 \\
    H3 & detector-targeted clarification & 0.537 & 0.537 & 0.162 & 2.48 & 4.43 \\
    H4 & sibling functions with doctests & 0.517 & 0.517 & $-0.032$ & 0.00 & 4.74 \\
    B2 & five direct samples instead of questions & 0.513 & 0.547 & $-0.065$ & 0.00 & 5.00 \\
    B3 & one-call baseline, wish and stub only & 0.310 & 0.310 & $-2.033$ & 0.00 & 1.00 \\
    H3$^-$ & H3's questions, answers withheld & 0.427 & 0.427 & $-0.903$ & 2.48 & 4.09 \\
    \bottomrule
  \end{tabular}
  \caption{Measured 100-task pilot under the corrected protocol, three samples per condition
  and task, $n=300$ units per condition. Recovery is the fraction of the B1--H0 gap closed;
  B2's pass rate uses all five internal samples. Source:
  \protect\path{results/runs/r2b-pilot/report.md} and
  \protect\path{results/runs/r2b-pilot/summary.json}.}
  \label{tab:conditions}
\end{table}

\begin{table}[!ht]
  \centering
  \small
  \begin{tabular}{@{}lrr@{}}
    \toprule
    Contrast & difference (pp) & 95\% CI (pp) \\
    \midrule
    B1 $-$ H0 & $+10.3$ & [$+3.0$, $+17.7$] \\
    H2 $-$ H0 & $-7.7$ & [$-14.3$, $-1.0$] \\
    H3 $-$ H0 & $+1.7$ & [$-5.7$, $+9.0$] \\
    H4 $-$ H0 & $-0.3$ & [$-6.0$, $+5.3$] \\
    H1 $-$ H0 & $-23.3$ & [$-31.3$, $-15.7$] \\
    H3 $-$ H2 & $+9.3$ & [$+1.7$, $+17.0$] \\
    H3 $-$ H4 & $+2.0$ & [$-6.3$, $+9.7$] \\
    H3$^-$ (no answers) $-$ H3 & $-11.0$ & [$-18.7$, $-2.7$] \\
    B2 $-$ H0 & $-0.7$ & [$-6.7$, $+5.0$] \\
    \bottomrule
  \end{tabular}
  \caption{Paired task-level bootstrap, 2{,}000 resamples, seed \texttt{20261121}, mean over
  each task's three pass@1 samples. Only B1--H0, H2--H0, H1--H0, H3--H2 and H3$^-$--H3
  exclude zero. Source: \protect\path{results/runs/r2b-pilot/bootstrap.json}.}
  \label{tab:bootstrap}
\end{table}

\subsection{Per-category analysis}

Table~\ref{tab:categories} groups tasks by an L0 field whose separately rendered content
their wish omits (the prose may still mention it). The pilot does not run the separate L2$_c$
single-field-drop conditions, so this grouping is attribution by task subset, not a controlled
one-field deletion; tasks can carry several categories, so the task counts overlap. H3 recovery is positive on all four subsets: constraints 0.194 (57 tasks),
signature and types 0.162 (all 100 tasks), examples 0.100 (99 tasks) and error behaviour 0.077
(10 tasks). H2 is below H0 on every subset, which is consistent with generic questions costing
turns without recovering the missing field. The error-behaviour subset is small (10 tasks,
30 units), so its estimate is exploratory. The examples row is confounded by the BigCodeBench
residual: 54 pilot tasks are BigCodeBench, whose assertion-based visible tests still expose
worked examples even when the examples field is deleted at L3.

The targeted condition's largest gain is on constraints, which is also the category the
detector ranks best (AUROC 0.952). That is the encouraging pattern: when the detector knows
which field is missing and the stakeholder can supply it, H3 recovers more than H0. The gain
is still partial: on constraints H3 passes 48.5\% against H0 45.0\% and B1 63.2\%, so most of
the subset's specification gap remains. On examples and error behaviour the subset estimates
are smaller, and the error subset cannot carry much weight.

\begin{table}[!ht]
  \centering
  \small
  \begin{tabular}{@{}lrrrrrrr@{}}
    \toprule
    Category & tasks & B1 & H0 & H2 & H3 & H4 & H3 recovery \\
    \midrule
    signature and types & 100 & 0.623 & 0.520 & 0.443 & 0.537 & 0.517 & 0.162 \\
    examples & 99 & 0.623 & 0.522 & 0.448 & 0.532 & 0.522 & 0.100 \\
    constraints & 57 & 0.632 & 0.450 & 0.380 & 0.485 & 0.450 & 0.194 \\
    error behaviour & 10 & 0.833 & 0.400 & 0.367 & 0.433 & 0.467 & 0.077 \\
    \bottomrule
  \end{tabular}
  \caption{Pass rates on the subset of pilot tasks whose L0 carries each category (the
  separately rendered field is absent at L3, though the prose may mention it). Categories
  overlap and the nine-condition pilot does not run separate L2$_c$ conditions. Source:
  \protect\path{results/runs/r2b-pilot/summary.json}, \texttt{by\_category}; the same
  grouping is mirrored in \protect\path{results/runs/r2b-pilot/bootstrap.json}.}
  \label{tab:categories}
\end{table}

\subsection{Detector results}

Table~\ref{tab:detector} gives five-fold out-of-fold detection on the 500 smoothed train
wishes. Intent and signature/types are positive on all 500 train tasks, so their AUROC is
undefined; the detector is effectively a three-label head with 293 constraints, 491 examples
and 85 error-behaviour positives. The encoder reaches macro-F1 0.792 (0.654 over the three
varying labels), AUROC 0.632 on examples, 0.952 on constraints and 0.744 on error behaviour,
with Brier 0.0439 and ECE 0.0142. TF-IDF is close on the aggregate (macro-F1 0.789, varying
0.648) and on examples/constraints (0.871/0.947), while the encoder wins on error behaviour
(0.744 vs 0.695); calibration is comparable (Brier 0.0470 vs 0.0439). The encoder's practical
advantage is that it scores a raw wish without a task-specific lexical model, not that it wins
every metric.

Error behaviour remains the weak label: 85 positives in the 500 train tasks, average
precision 0.315, and F1 at the deployed 0.5 threshold is 0.046 because calibrated
probabilities stay near the 17\% base rate. AUROC and average precision are the informative
ranking metrics for that label. Examples are nearly constant (491/500 positive, prevalence
0.982), so its AUROC is computed from only 9 negatives and is not informative; F1 0.991 mostly
reflects the base rate. The H3 policy selects constraints on 70\% of tasks and
error behaviour on 37\%; intent is never selected.

\begin{table}[!ht]
  \centering
  \small
  \setlength{\tabcolsep}{3.5pt}
  \begin{tabular}{@{}lccccccc@{}}
    \toprule
    Method & macro-F1 & macro-F1 (varying) & AUROC ex. & AUROC cons. & AUROC err. & Brier & ECE \\
    \midrule
    Prior & 0.746 & 0.577 & 0.251 & 0.442 & 0.472 & 0.0806 & 0.0096 \\
    TF-IDF & 0.789 & 0.648 & 0.871 & 0.947 & 0.695 & 0.0470 & 0.0131 \\
    Encoder & 0.792 & 0.654 & 0.632 & 0.952 & 0.744 & 0.0439 & 0.0142 \\
    \bottomrule
  \end{tabular}
  \caption{Five-fold out-of-fold detection on the 500 smoothed train tasks. ``Varying'' averages
  examples, constraints and error behaviour; intent and signature/types are positive on every
  task, so their AUROC is undefined. Source:
  \protect\path{results/runs/r2b-pilot/detector-smoothed/metrics.json}.}
  \label{tab:detector}
\end{table}

\subsection{Failure analysis and interpretation}

Across all 2{,}700 units there are 1{,}263 passes and 1{,}437 failures under the any-sample
rule used for pass rate (1{,}253 under first-sample pass@1). The failure buckets are dominated
by \texttt{AssertionError} (858, 59.7\%), followed by \texttt{NoSolution} (281,
19.6\%), \texttt{TypeError} (92, 6.4\%), \texttt{ValueError} (63, 4.4\%) and
\texttt{AttributeError} (61, 4.2\%). The runner's \texttt{NoSolution} bucket means the agent
never wrote \texttt{solution.py} and did not pass; the assertion bucket is a wrong
implementation that ran. This distribution is what a hidden-suite pilot should look like: most
failures are semantic, not infrastructure. Only 26 units time out and 8 end in
\texttt{XMLSyntaxError}.

The corrected negative and positive results should be read together. Targeting is doing
something: H3 beats H2 and recovers 16\% of the measured gap, and H3$^-$ is clearly below H3,
so the dialogue channel matters. It is not doing enough: H3 is within 1.7 points of H0 and
2.0 points of H4, and both intervals cross zero. Three mechanisms are plausible. First, the
categories the detector ranks best may be the ones the model can already partly infer from the
stub and visible tests, so the marginal value of an answer is limited. Second, H3 compiles its
answers into a Requirements block and one repair call, while H2 can keep asking and iterating
for up to five questions; the generic condition's 6.36 turns are not buying accuracy and may
be adding noise. Third, the weakest category, error behaviour, has only 10 tasks in the
pilot, so the subset that most needs targeted clarification is the least powered. The full
study should separate these mechanisms with L2$_c$ conditions, a larger error-behaviour
sample, and the parser/visible-test fixes now in place.

\section{Limitations}
\label{sec:limitations}

\textbf{Pilot size and power.} The pilot uses the 100-task development split, not the 300
held-out tasks. It has one seed base and three samples per condition; the intervals describe
task-level uncertainty only, not model-seed or session variance. A 10-point gap is detectable,
but a 1--2-point difference among H2, H3 and H4 is not, and the per-category error-behaviour
subset (10 tasks) is exploratory.

\textbf{One session, transport and GPU.} All units come from one A100-SXM4-40GB session with
vLLM 0.34.0 and the JSON text protocol because native tool calls were unavailable on that
build. Results may differ under native tool calling, a different vLLM build, a larger GPU or
multiple sessions.

\textbf{Compute and budget.} The run used 2.902 GPU-hours on one A100-SXM4-40GB and \$0 paid
API. The pilot is therefore a measured development run, not the final evaluation; the full
study remains future work.

\textbf{Benchmark composition and the BigCodeBench residual.} BigCodeBench supplies 54 of the
100 pilot tasks. HumanEval+ (12) and MBPP+ (34) are small and may not represent the library
call stress case. The pool is 1{,}100 kept tasks because of the fixed split and validation
exclusions; the held-out 300 remain unrun. BigCodeBench visible tests remain assertion-based,
so the examples category is not fully deleted there and the examples row is confounded.

\textbf{Corrected protocol and H1.} The whole-prose intent keeps later requirements in the
wish, which is the intended correction but also means the L3 wish is less degraded than the
name suggests. Inputs-only visible tests remove expected values for HumanEval+/MBPP+, which
makes the H1 acceptance-test condition weak by construction: its 28.7\% is partly a protocol
artifact, not a fair test of acceptance-test-first practice. The smoothing step falls back to
the raw intent for 12 of 600 wishes.

\textbf{Detector and levels.} The detector is a frozen encoder with a linear probe, not a
fine-tuned model, and the committed head uses 600 full-batch AdamW steps (learning rate 0.01,
L2 0.3) rather than the proposal's three-epoch minibatch recipe; its aggregate edge over TF-IDF
is small, and its error-behaviour F1 at 0.5 is low. The pilot does not run the L2$_c$
single-field-drop conditions, so per-category attribution is by task subset. H5 is not
implemented.

\textbf{Recovery denominator.} Recovery uses B1 and H0 measured on the same 100 tasks, so the
denominator is itself noisy. We report paired task-level intervals for contrasts and treat
recovery as a descriptive normalization, not a precise effect size.

\section{Conclusion and future work}

The corrected measured pilot supports one strong claim and one relative claim. The
specification gap is real: removing the full specification costs 10.3 points, and that
interval excludes zero. Targeted clarification beats generic clarification by 9.3 points and
is the only condition that recovers any of the gap at the point estimate (recovery 0.16). That
is a positive result for detector-targeted elicitation, where the earlier measurement found a
null under a parser that deleted later prose from L3 and visible tests that restored expected
values. But targeting does not beat the ceiling or sibling-function
context, its differences from H0 and H4 cross zero, generic questions hurt, withholding the
targeted answers hurts more, and acceptance tests first is catastrophic under inputs-only
visible tests. The detector ranks constraints well but its error-behaviour label remains
weak.

The full study should scale to the 300 held-out tasks with multiple seeds and sessions, add
the L2$_c$ single-field conditions with examples truly removed, enlarge the error-behaviour
sample, and test H5. It should also settle whether the generic-question penalty is a real
property of generic dialogue or a turn-budget artifact by matching turns across clarification
conditions. Until then, the defensible statement is relative: targeted questioning is better
than generic questioning on this 100-task pilot, and the specification gap remains open.

\section{Reproducibility and artifact map}
\label{sec:reproduce}

All empirical numbers in this report come from the corrected run under
\path{results/runs/r2b-pilot/}:

\begin{itemize}
  \item \path{summary.json}, \path{report.md} and \path{results.jsonl}: 2{,}700 unit
        records, condition aggregates, recovery, per-category grouping and failure buckets;
  \item \path{bootstrap.json} and \path{bootstrap.md}: the paired task-level contrasts in
        Table~\ref{tab:bootstrap} and the per-category grouping in
        Table~\ref{tab:categories}, produced by \path{src/wishrec/harness/bootstrap.py};
  \item \path{manifest.json}, \path{environment.json}, \path{endpoint.json},
        \path{gpu_time.json} and \path{data_provenance.json}: provenance, environment,
        transport probe, shard timing and corrected-levels hashes;
  \item \path{detector-smoothed/}: the smoothed five-label refit and its metrics;
  \item \path{smoothing/}: the 600 smoothed wishes (500 train, 100 pilot; 12 fallbacks) and
        their manifest;
  \item \path{units/}: the per-unit traces; \path{dry-run/}: the 45-unit plumbing check;
  \item \path{COLAB_RUN.md} and \path{colab_a100_vllm.ipynb}: the operator's runbook and
        executed notebook.
\end{itemize}

Data counts, split sizes and the corrected protocol come from \path{results/levels/stats.json}
and \path{results/levels/manifest.json}; the detector artifact is under
\path{results/detector-smoothed/}. The originally submitted milestone is preserved under
\path{paper/milestone/submitted/}. The paper is compiled with pdf\LaTeX{} and BibTeX from
\path{paper/main.tex}.

\bibliographystyle{unsrt}
\bibliography{references}

\end{document}
